\chapter{Metodología: Spec-Driven Development}\label{cap:Metodologia}
% Índice v4, cap. 3: el proceso de forma general, sin resultados (van en el cap. 4).

\section{Principios}\label{sec:Principios}
% SDD spec-anchored (acordado con el tutor): la especificación es la fuente de
% verdad y se mantiene viva; no se promete regenerar código idéntico; la
% conformidad la garantizan las pruebas y /speckit-converge.
% Separación qué/cómo; cambios primero en la especificación; humano en el bucle;
% criterios de aceptación convertidos en pruebas antes de implementar.
\pendiente{principios de SDD spec-anchored}

\section{Artefactos}\label{sec:Artefactos}
% Tabla: constitution.md, spec.md, plan.md, tasks.md y CLAUDE.md, con a qué
% responde cada uno, quién lo produce y quién lo aprueba.
\pendiente{tabla de artefactos}

\section{Fases del ciclo}\label{sec:Fases}
% Siete fases; la implementación es un bucle (loop engineering) con límite de
% iteraciones y /speckit-converge como criterio de parada (acordado con el tutor).
% Cada funcionalidad recorre las fases 2-7 en una rama corta.
\pendiente{fases del ciclo y diagrama}

\section{Trazabilidad y registro del proceso}\label{sec:Trazabilidad}
% Base: docs/conceptos/03-git-github-y-trazabilidad.md
% Cadena REQ-xx -> T-xx -> commit -> prueba; registro de prompts e iteraciones;
% métricas del proceso.
\pendiente{trazabilidad y métricas del proceso}

\section{Riesgos del desarrollo con IA y su control}\label{sec:Riesgos}
% Alucinaciones, dependencias inseguras, secretos y deuda técnica; revisión
% humana, análisis estático y de dependencias, pruebas obligatorias y límites.
\pendiente{riesgos y medidas de control}

\section{Planificación temporal}\label{sec:Planificacion}
% Iteraciones por funcionalidad en un diagrama de Gantt.
\pendiente{diagrama de Gantt}
